\chapter{Conclusiones y trabajo futuro}\label{cap:conclusiones}

\section{Conclusiones}

Este Trabajo Fin de Grado en Ingeniería Informática ha cubierto el
diseño y la implementación de una plataforma SaaS completa para la
detección y verificación automática de noticias falsas, con cuatro
componentes desplegables independientes (backend, frontend, extensión
y base de datos) y un agente de verificación basado en FEVER expuesto
como funcionalidad diferencial del plan \emph{Super Pro}.

Las contribuciones más relevantes son:

\begin{enumerate}[leftmargin=*,itemsep=0.2em]
  \item Una arquitectura desacoplada con flujo unidireccional de
        datos en el frontend y separación por dominio en el backend.
  \item La integración real con Supabase Auth, RLS, Stripe Billing y
        una API de búsqueda externa para evidencias.
  \item Una extensión MV3 con renderizado seguro y reutilización de
        autenticación y servicios.
  \item Una estrategia de pruebas automatizada que cubre los módulos
        críticos.
  \item Un despliegue contenedorizado reproducible.
\end{enumerate}

\section{Trabajo futuro}

\begin{itemize}[leftmargin=*,itemsep=0.2em]
  \item Sustituir la API de búsqueda externa por un recuperador propio
        sobre un índice mantenido (Elasticsearch / OpenSearch).
  \item Añadir soporte multilenguaje real al modelo (modelos NLI
        multilingües).
  \item Migrar el despliegue a Kubernetes con \emph{autoscaling}
        horizontal.
  \item Incorporar OpenTelemetry para \emph{tracing} y métricas
        Prometheus.
  \item Publicar la extensión en Chrome Web Store y Firefox Add-ons
        tras una auditoría de privacidad.
\end{itemize}
